% Wave-14 Q1/20 --- slides Jacobi / $\phi_{0,1}$ / $\Lambda^{3,2}$ extraction, lines 100--200
% Primary source: /tmp/automorphic-slides.txt (Lorgat 2020 slides)
% Cross-check: Lorgat 2020 paper \S 2--3, Eichler--Zagier 1985, Gritsenko--Nikulin 1998

\providecommand{\kBKM}{\kappa_{\mathrm{BKM}}}
\providecommand{\kch}{\kappa_{\mathrm{ch}}}
\providecommand{\Sp}{\mathrm{Sp}}
\providecommand{\GL}{\mathrm{GL}}
\providecommand{\SL}{\mathrm{SL}}
\providecommand{\SO}{\mathrm{SO}}
\providecommand{\Orth}{\mathrm{O}}
\providecommand{\HH}{\mathbb{H}}
\providecommand{\Zd}{\mathbb{Z}}
\providecommand{\Cd}{\mathbb{C}}

\section*{Wave-14 Q1/20: slides Jacobi / $\Lambda^{3,2}$ extraction (lines 100--200)}

\subsection*{Cycle 1 --- Jacobi-form data on display}

Slide lines 100--177 unfold the self-structure of Igusa's $\Delta_5$ at
the Fourier-Jacobi boundary. Four data points are displayed:

\emph{Coefficient normalisation.} $f(1,1,1)=64$ and $64\mid f(n,l,m)$
for every Fourier coefficient of $\Delta_5(Z)=\sum f(n,l,m)
\exp(\pi i(nz_1+lz_2+mz_3))$ over the summation locus $n,l,m\equiv 1
\pmod{2}$, $4nm-l^2>0$, $n,m>0$ (lines 96--103, 111). The factor $64$ is
not decorative: it is the Gram determinant needed to normalise the
first Fourier--Jacobi coefficient against $\eta^{18}\nu_{11}^2$.

\emph{Macdonald-type identity} (line 115):
\[
1+\frac{1}{64}\sum_{t\in\mathbb{N}}f(1+2t,1,1)q^t
=\prod_{k\in\mathbb{N}}(1-q^k)^9.
\]

\emph{Fourier--Jacobi decomposition} (line 120):
\[
\Delta_5(Z)=\sum_{\substack{m>0\\m\equiv 1\,(2)}}
\phi_{5,m}(z_1,z_2)\exp(\pi i m z_3),
\]
with $\phi_{5,m}$ a Jacobi cusp form of index $m/2$ and non-trivial
Maass multiplier character.

\emph{Jacobi theta-series and triple-product} (lines 129--145):
\[
\nu_{11}(z_1,z_2)=\sum_{n\in\mathbb{Z}}(-1)^n
\exp\bigl(\tfrac{\pi i}{4}(2n+1)^2 z_1+\pi i(2n+1)z_2\bigr),
\]
\[
\nu_{11}=q_1^{1/8}q_2^{-1/2}\prod_{n\geq 1}
(1-q_1^{n-1}q_2)(1-q_1^n q_2^{-1})(1-q_1^n).
\]
The assembly $\psi_{5,1/2}=\eta(z_1)^9\,\nu_{11}(z_1,z_2)$ (line 147) is
the unique-up-to-scale Jacobi cusp form of weight $5$ and index $1/2$
with the $\Delta_5$ character; squaring lands in weight~$10$ index~$1$
and matches the first Fourier--Jacobi coefficient of $\Delta_{10}=
\Delta_5^2$:
\[
\tfrac{1}{64}\phi_{5,1}=\psi_{5,1/2}^2
=-q_1^{1/2}q_2^{-1}\prod_{n\geq 1}(1-q_1^{n-1}q_2)(1-q_1^n q_2^{-1})
(1-q_1^n)^{10}.
\]
The Macdonald identity of line 115 is a corollary by extracting the
coefficient of $q_2^{1/2}$ and applying the Jacobi triple-product
(lines 174--177).

\subsection*{Cycle 2 --- $\phi_{0,1}$ against Eichler--Zagier 1985}

The slides' 100--200 window does \emph{not} yet display $\phi_{0,1}$
itself; that weak Jacobi form of weight~$0$ index~$1$ enters later (\S 6
of the paper). At this stage the slides are preparing the Borcherds-lift
half-space: the explicit product formula for $\psi_{5,1/2}^2$ establishes
the building block whose Borcherds lift is $\Delta_5$. The
Eichler--Zagier theory (Eichler--Zagier, \emph{The Theory of Jacobi
Forms}, Progr.\ Math.~55, Birkh\"auser 1985) supplies the
weight-$0$ index-$1$ weak Jacobi form
\[
\phi_{0,1}(z_1,z_2)=\phi_{12,1}(z_1,z_2)/\delta_{12}(z_1)
=(r^{-1}+10+r)+q(10r^{-2}-64r^{-1}+108-64r+10r^2)+\cdots
\]
(paper \S 6, p.~9), with $\delta_{12}(z_1)=q\prod_{n\geq 1}(1-q^n)^{24}$
and $\phi_{12,1}$ the first Fourier--Jacobi coefficient of
$\Delta_{12}$. The leading term $r^{-1}+10+r$ matches
Eichler--Zagier's classical generator of the weak-Jacobi ring at
weight~$0$: the $10$ is the codimension one would expect from the
$E_{4,1}/E_{4,0}$ decomposition, and the expansion is consistent with
the Eichler--Zagier tables. No discrepancy between slide-stage prep and
Eichler--Zagier.

\subsection*{Cycle 3 --- $\Lambda^{3,2}$ against Gritsenko--Nikulin 1998}

Slide lines 178--239 execute the $\wedge^2$ transfer. The lattice data:
\[
\Lambda^4=\Zd e_1\oplus\Zd e_2\oplus\Zd e_3\oplus\Zd e_4,\qquad
\wedge^2:\Lambda^4\to\Lambda^4\wedge\Lambda^4,
\]
with $\SL_4(\Zd)$-invariant Pfaffian pairing
$(,):\Lambda^4\wedge\Lambda^4\to\Cd$, $u\wedge v=(u,v)e_1\wedge\cdots
\wedge e_4$, of signature $(3,3)$. Fixing $q=e_1\wedge e_3+e_2\wedge
e_4$ picks out
\[
\{g:\Lambda^4\to\Lambda^4\mid(g\wedge g)q=q\}\simeq\Sp_4(\Zd),
\]
and
$\Lambda^{3,2}:=q^{\perp}\subset\Lambda^4\wedge\Lambda^4$ is even
unimodular-up-to-$[2]$-summand of signature $(3,2)$:
\[
\Lambda^{3,2}\simeq\Lambda^{(1,1)}\oplus\Lambda^{(1,1)}\oplus[2].
\]
The basis $(f_1,f_2,f_3,f_{-2},f_{-1})=
(e_1\wedge e_2,\,e_2\wedge e_3,\,e_1\wedge e_3-e_2\wedge e_4,\,
e_4\wedge e_1,\,e_4\wedge e_3)$ (line 219) diagonalises the Gram into
$2U\oplus\langle 2\rangle$. This is precisely the paramodular
$\Sp_4\mid\Orth(2U\oplus\langle 2t\rangle)$ transfer at $t=1$ of
Gritsenko--Nikulin 1998 (``Siegel automorphic form corrections of some
Lorentzian Kac--Moody Lie algebras,'' Amer.\ J.\ Math.~119), and the
subsequent type-IV domain $\HH^{IV}=\{Z\in\mathbb{P}(\Lambda^{3,2}
\otimes\Cd)\mid(Z,Z)=0,(Z,\bar Z)<0\}$ with the identification
$\HH^{IV}_+\simeq\HH_2$ under $Z=\binom{z_2\ z_3}{z_1\ z_2}$ (line
238) is the standard GN Type~IV realisation. Consistency: exact match.

\subsection*{Cycle 4 --- what Wave 13 missed}

Wave 13 P1 (\texttt{notes/wave13\_p1\_slides\_intro\_context.tex})
enumerated items 3, 7, 8 of the game-plan at outline level. The 100--200
window adds four load-bearing pieces not in P1:

(a) The explicit isomorphism
$\wedge^2:\Sp_4(\Zd)/\{\pm I_4\}\xrightarrow{\sim}\SO_+(\Lambda^{3,2})
\simeq\Orth(\Lambda^{3,2})_+/\{\pm I_5\}$, with the Pfaffian pairing as
its intertwiner.

(b) The theta-decomposition $\psi_{5,1/2}=\eta^9\nu_{11}$, which is
the explicit Jacobi-theta factorisation underlying the $\phi_{5,1}$
Fourier coefficients and hence the Borcherds denominator.

(c) The $64\mid f(n,l,m)$ divisibility (line 111), a Gram-level
consistency condition not visible from the outline.

(d) The Macdonald-type identity of line 115 as a direct corollary of
the Jacobi triple-product on $\psi_{5,1/2}^2$; this is the explicit
bridge between $\Delta_5$'s Fourier coefficients and the $\prod(1-q^k)
^9$ side that later becomes the Weyl--Kac denominator of
$\mathfrak{g}_{\Delta_5}$.

Each of (a)--(d) is cited in paper \S 2--3 and must be referenced in
any $\kBKM(\Phi_1)=5$ exposition in Vol III.

\subsection*{Cycle 5 --- frontier: game-plan step at this stage}

The slide window closes (line 207--239) mid-execution of game-plan
item~3 (``trivial transfer $\Sp_4(\Zd)\to\Orth(\Lambda^{3,2})^+$'').
Items~1--2 are complete (diagonal-divisor classification; $\Delta_5$
self-structure via Fourier--Jacobi + $\nu_{11}$ + triple product).
Item~3 is in progress: the $\wedge^2$ isomorphism and the $\HH^{IV}
\simeq\HH_2$ identification are displayed, the connected-component
count $\Orth_{\mathbb{R}}(\Lambda^{3,2})$ (four components, of which
two preserve $\HH^{IV}_+$) is pending at line 239. Items~4--9 (reflection
groups $W_I^{(2)}, W_{II}^{(2)}$; Weyl vector $\rho$; hyperbolic
sublattice $\Lambda^{2,1}_{II}$; automorphic correction
$\mathfrak{g}_{\Delta_5}$; root multiplicities via $\phi_{0,1}$; CY3
arithmetic) remain for later slides. The $\kBKM(\Phi_1)=c_1(0)/2=5$
payoff lands at item~8.
